\documentclass[12pt]{article}
\usepackage{../../includes/lecture_notes}
\usepackage{../../includes/math}
\usepackage{../../includes/uark_colors}

\title{Forecasting}
\author{Prof. Kyle Butts}
\date{}

\hypersetup{
  colorlinks = true,
  allcolors = ozark_mountains,
  breaklinks = true,
  bookmarksopen = true
}

\begin{document}
\maketitle
\begin{multicols}{2}
  \section*{Introduction to Forecasting}

  The goal of forecasting is to learn the relationship between \emph{input variables} and \emph{outcome variable(s)} so that we can \emph{predict} the outcome when we don't observe it.

  This shows up everywhere. Say we want to learn who are potential customers based on their observable characteristics:
  \begin{itemize}
    \item[] \textbf{Input:} observable characteristics (age, income, location)
    \item[] \textbf{Outcome:} whether they purchase a product
  \end{itemize}

  Or we might want to predict values of a variable in the future, like the time-series of stock prices:
  \begin{itemize}
    \item[] \textbf{Input:} the time-period
    \item[] \textbf{Output:} stock price
  \end{itemize}

  In both cases, we're trying to learn some relationship between the inputs and the outputs so we can make predictions when we don't have the outcome.

  \section*{The Forecasting Model}

  We have an outcome variable $y$ and a set of $p$ different predictor variables $\bm{X} = (X_1, X_2, \dots, X_p)$.
  For some observations we observe both $\bm{X}$ and $y$.
  This is what allows us to \text{fit} our model and learn the relationship between the two.

  We can write our model in general form as
  $$y = f_0(\bm{X}) + \varepsilon,$$
  where $f_0$ is some unknown function of $\bm{X}$ and $\varepsilon$ is the \emph{error term}.

  Think of $f_0(\bm{X})$ as the `true' model: this is the best prediction of $y$ given information on $\bm{X}$.
  We can write this formally as $\expec{y - f_0(\bm{X})}{\bm{X}} = \expec{\varepsilon}{\bm{X}} = 0$.
  This is a really strong assumption.
  It says that we have \emph{exhausted} all the information in $\bm{X}$.
  Any of the remaining variation in $y$ is completely unpredictable given the information we observe in $\bm{X}$, even if we had access to the best prediction model.
  % Once we know $f_0(\bm{X})$, we cannot predict any more of the variation in $y$ using $\bm{X}$.

  \section*{Predictions and the Error Term}

  We will use a training sample of data to predict $f_0$.
  We denote our estimate of $f_0$ as $\hat{f}$.
  For a given unit, we predict
  $$\hat{y} = \hat{f}(\bm{X}).$$

  Definitionally, we have a \textbf{prediction error} of $\hat{f}$ which is defined as $y - \hat{f}(\bm{X})$.
  The term `error'' gets used in a few different ways, so let me clarify:

  \begin{enumerate}
    \item If we knew the `true' model, $f_0(\bm{X})$, then $\varepsilon \equiv y - f_0(\bm{X})$ represents the things that are unpredictable given $\bm{X}$.
      This is the ``true'' error term.

      \smallskip
    \item If $\hat{f}(\bm{X})$ is an estimated model, then $\hat{\varepsilon} = y - \hat{f}(\bm{X})$ is the difference between that unit's actual $y_i$ and their predicted $y$.
      This is better called the \emph{prediction error}.
  \end{enumerate}

  The key distinction: the true error term is a property of the best possible model $f_0$ while the prediction error is a property of our estimated model, $\hat{f}$

  \section*{Goals of Forecasting}

  There are two related goals when predicting $f$:

  \noindent\textbf{1. Prediction:} Predict $y$ as well as possible.
  \begin{itemize}
    \item Think of prediction as a `black box' where the goal is to do as good a job at predicting $y$ as possible.
    \item We want to learn the `true' $f_0(\bm{X})$, leaving no information on the table.
  \end{itemize}

  \noindent\textbf{2. Inference:} Understand the relationship between $\bm{X}$ and $y$.
  \begin{itemize}
    \item If our goal is being able to \emph{describe} the relationship between $x$ and $y$; i.e. we care about understanding $f_0$ and not just $\hat{y}$.

    \item Allows us to create ``heuristics'' that help decision-makers.

    \item Worth approximating $f_0$ with a more `simple' functional form so that we can better convey the results to stakeholders.
  \end{itemize}

  For example, with housing data:
  \begin{itemize}
    \item[] \textbf{Prediction example:} ``Is the house over/under-priced?''. Only care about $\hat{\text{price}} = \hat{f}(\bm{X})$.
    \item[] \textbf{Inference example:} ``How much more do homes with a river view sell for?''.  Need to know information about $f_0$ to answer.
  \end{itemize}

  \section*{Model Flexibility}

  There is a tension between how \emph{flexible} we can make our model and how well it performs:

  \begin{enumerate}
    \item If our goal is prediction, we only have a finite amount of data to use to fit our model, so there's a limit on how much we can learn. We face the risk of \textbf{overfitting} the data.
      Overfitting means that we are chasing after the random noise $\varepsilon$ instead of the true signal.

    \item If our goal is inference, then added flexibility makes our results harder to summarize and communicate to stakeholders.
  \end{enumerate}

  This is the core trade-off in forecasting: more flexible models can capture more complex patterns, but they also introduce more noise and are harder to interpret.

  \section*{Fitting Models}

  The general framework for forecasting follows these steps:

  \begin{enumerate}
    \item Collect a set of data $\{ (y_i, \bm{X}_i) \}_{i=1}^n$ for a set of observations.

      \smallskip
    \item Using knowledge of the topic, select a \emph{class} of models $\mathcal{F}$ that you want to select from.
      \begin{itemize}
        \item That is, $\hat{f}$ must be one of the functions in $\mathcal{F}$, e.g. linear functions of $\bm{X}_i$.
      \end{itemize}

      \smallskip
    \item Using your data, select $\hat{f} \in \mathcal{F}$ that \emph{predicts} $y$ ``best''.
      \begin{itemize}
        \item ``best'' is defined by a \emph{loss function}, e.g. mean-squared prediction error.
      \end{itemize}
  \end{enumerate}

  \section*{Selecting a Class of Models}

  First, we have to select a class of models $\mathcal{F}$. This involves:
  \begin{itemize}
    \item Selecting which variables we want to include in our model.
    \item Specifying a functional form for our model.
  \end{itemize}

  For example, we might think about a simple \emph{linear model}:
  $$f(\bm{X}_i) = \alpha + X_{i1} \beta_1 + X_{i2} \beta_2.$$

  The class $\mathcal{F}$ consists of all models of this form, i.e. choices of different values of $(\alpha, \beta_1, \beta_2)$.
  This is a restrictive model:
  \begin{itemize}
    \item No polynomials of $X_1$ (e.g., wages are non-linear in age).
    \item No interaction between $X_1$ and $X_2$ (e.g., a college degree changes the return to experience).
  \end{itemize}

  Perhaps we want a wider class of models:
  $$\alpha + X_{i1} \beta_1 + X_{i1}^2 \beta_2 + X_{i2} \beta_3 + X_{i2}^2 \beta_4 + X_{i1} X_{i2} \beta_5.$$

  The class $\mathcal{F}$ consists of all functions of this form, and since some coefficients can be zero, this class is \emph{strictly more general} than the last. We can imagine creating even more terms or using things other than quadratics, e.g., $\one{20 < \text{Age}_i \leq 30}$.

  The choice of $\mathcal{F}$ is the `art' of being a good forecaster.
  We want to include functions of $\bm{X}$ that we think are really useful for predicting $y$ and the only way to know these functions is to become a subject matter expert on whatever phenomenon you are modelling.

  Later in your studies, you will learn about machine learning algorithms that ``learn'' the transformations of $\bm{X}$ for you.
  This sounds like a `free lunch' type scenario, but remember that you will only have a finite amount of data.
  Telling the model important transformations to include will result in more precise forecasts than making the machine learning method learn them for you.


  \section*{Loss Functions}

  To provide a summary measure of fit, we want to \emph{average} prediction error over many observations.
  If we took the simple mean of prediction error, positive and negative errors would cancel out.
  For example, an error of $-1$ and $1$ would be just as bad as $-4$ and $4$.

  \bigskip
  Therefore, we square the errors so that they are all positive and do not cancel out.
  The most common loss function is the \emph{mean-squared (prediction) error} (MSE):
  $$
  \text{MSE} \equiv \frac{1}{n} \sum_{i=1}^n \left( y_i - \hat{y}_i \right)^2 = \frac{1}{n} \sum_{i=1}^n \hat{\varepsilon}_i^2.
  $$

  The MSE is not the only loss function. Consider:
  \begin{itemize}
    \item The \emph{mean-absolute prediction error} $$\frac{1}{n} \sum_{i=1}^n \left| y_i - \hat{y}_i \right|$$ will estimate the \emph{median} of $y$ given $\bm{X}$.
    \item In settings where you're predicting whether someone has a disease, you might want to penalize false-negatives more than false-positives.
  \end{itemize}

  \section*{In-Sample vs. Out-of-Sample Prediction}

  As a forecaster, you will \emph{fit} a model using a set of observations $\{ (y_i, \bm{X}_i) \}_{i=1}^n$. This is called the \emph{training data}.

  We can calculate the \emph{in-sample MSE} by averaging over all observations in the training data. This tells us how well we predict the data we trained our model on.

  But if our goal is prediction, we really want to know how our model would predict on \emph{new} observations we haven't seen before.
  It is common to hold out a set of \emph{test data} that is NOT used for training our model, but only for evaluating its performance.

  It is tempting to pick our model with the smallest \emph{in-sample MSE}, but this can be \emph{a risky thing to do}.
  By focusing on fitting the current sample very well, you risk \emph{overfitting} the data.
  Your model starts chasing the random noise $\varepsilon$ instead of learning the true pattern.

  By making our model more and more flexible, you risk overfitting more and more.
  Your model tries to improve the \emph{in-sample} MSE, but worsens your \emph{out-of-sample} MSE.
  The solution is to evaluate your model using held-out test data.

  \section*{Sample Splitting}

  One solution to this problem is \textbf{sample splitting}.
  Say you want to fit a polynomial but are concerned with overfitting.
  We can tackle this with sample splitting:
  \begin{itemize}
    \item Take a random half your data and fit a polynomial of order $k$.
    \item Evaluate MSE on the other half (the test set).
  \end{itemize}

  Do this for $k = 1, \dots, K$ and pick the polynomial degree that minimizes test-data MSE.
  This technique is not very common for econometricians.
  I think the reason is that we tend to be concerned with \emph{inference} and therefore use relatively simple models (e.g., a regression with a few terms).
  Simplicity of the model class already provides some protection against overfitting, so sample splitting is less commonly used.

  %   \section*{The Bias-Variance Tradeoff}
  %
  %   The discussion of increasing flexibility leading to increasing noise in our model fit is a well-known problem. It is called the \emph{bias-variance tradeoff}:
  %
  %   \bigskip
  %   \begin{enumerate}
  %     \item \textbf{Bias:} When our model we fit, $\hat{f}(\bm{x})$, does a poor job fitting the true model $f_0(\bm{x})$.
  %     \item \textbf{Variance:} The variability of our model we fit, $\hat{f}(\bm{x})$, across different samples.
  %   \end{enumerate}
  %
  %   \bigskip
  %   This is a trade-off. To lower bias by adding flexibility, you typically add variance (noisiness) to the estimate. More flexible models can capture more complex patterns, but they vary more across samples.

  \section*{Inference on forecasts}

  Here's the crucial point that often gets glossed over: \emph{a forecast $\hat{y}$ is a function of the observed data.}

  Think about what this means. Your training data $\{ (y_i, \bm{X}_i) \}_{i=1}^n$ is a random sample from the population. Since $\hat{y} = \hat{f}(\bm{X})$ depends on this random data, $\hat{y}$ is itself a random variable.

  This is actually a profound observation.
  It is easy to think of predictions as just a number.
  But because predictions depend on randomly sampled data, they inherit the randomness from our sampling process.

  \section*{The Sample Distribution of a Forecast}

  Because $\hat{y}$ is a random variable (it's a function of random data), it has a \emph{sample distribution}. This is the distribution of $\hat{y}$ values you would get if you drew many different samples from the population and computed $\hat{y}$ each time.

  This concept is exactly the same as the sampling distribution of a sample mean or any other estimator.
  Just as $\bar{X}$ varies across samples, so does $\hat{y}$.

  Understanding the sample distribution will allow us to talk about the \emph{variability} of a forecast.
  Different samples yield different forecasts, and understanding this variability is crucial for understanding how reliable our predictions are.

  \section*{The Central Limit Theorem and Forecasts}

  Just like sample means, forecasts often have a sample distribution that is approximately \emph{normal}.
  This follows from a \emph{Central Limit Theorem} (CLT).
  The quality of the normal approximation will improve with the sample size.
  This is not always true (e.g. a complicated estimation procedure with forknig paths), so beware!

  The CLT tells us that, under relatively mild conditions, when we sum or average a large number of random variables (each of which might come from any distribution), the result tends toward a normal distribution.

  Since our forecast $\hat{y}$ is a function of many random data points, the CLT often applies, and the sample distribution of $\hat{y}$ will be approximately normal. This is true even if the underlying data or the true function $f_0$ are not normally distributed.

  This approximate normality is incredibly useful because the normal distribution has well-known properties that let us quantify uncertainty.

  \section*{Standard Errors and Confidence Intervals}

  Assuming $\hat{y}$ has an (approximately) normal sample distribution, we can use the tools of statistical inference to characterize our uncertainty.

  First, we need the \emph{standard error} of the forecast: the estimated standard deviation of the sample distribution of $\hat{y}$. This tells us how much $\hat{y}$ varies across repeated samples.

  Then, because the sample distribution is approximately normal, we know that approximately 95\% of $\hat{y}$ values across repeated samples fall within $\pm 1.96$ standard deviations of the true population value $f_0(\bm{X})$.

  This lets us form a \emph{confidence interval} for our forecast:
  $$\hat{y} \pm 1.96 \times \text{SE}(\hat{y}).$$
  The 95\% \textbf{coverage} of the confidence interval means that if we were to draw many samples and form a confidence interval for each sample, about 95\% of those intervals would contain the true $f_0(\bm{X})$.

  \section*{Putting It All Together}

  Let's trace through the logic:
  \begin{enumerate}
    \item A forecast $\hat{y}$ is a function of the randomly sampled training data.
    \item Therefore, $\hat{y}$ has a sample distribution
    \item The Central Limit Theorem typically ensures this sample distribution is approximately normal.
    \item We can estimate the standard deviation of this sample distribution (the standard error).
    \item Using the normal distribution, we can form confidence intervals around our forecast.
  \end{enumerate}

  This is the same logical structure as estimation in any statistical setting. Whether we're estimating a mean, a regression coefficient, or a forecast, the repeated-sampling perspective gives us a way to quantify statistical uncertainty.
\end{multicols}
\end{document}
